% Complete replacement for the template Abstract and Introduction.
% The abstract contains fewer than 300 words, excluding LaTeX commands.
% Citation keys already present in refs.bib are listed in the accompanying
% references_to_add.bib file; only genuinely new entries need to be merged.

\begin{abstract}
Earthquake declustering separates events judged to be background or sequence
mainshocks from events associated with earlier or later members of an
earthquake sequence. Different declustering algorithms can produce different
classifications. In a non-stationary epidemic-type aftershock sequence (ETAS)
model, these classifications may be used only to initialise a time-varying
background rate, but it is unclear whether their differences persist after
likelihood optimisation. This study assessed that sensitivity using 15,207
earthquakes of magnitude at least 2.0 from a Southern California catalogue
covering 20 May 2000 to 20 May 2026. Gardner--Knopoff, nearest-neighbour, and
Reasenberg-style declustering were applied to the same events. Their outputs
were converted into smooth background-rate starting functions, after which the
same penalised non-stationary temporal ETAS model was fitted from each start.

The methods classified 17.1\%, 30.1\%, and 25.3\% of events, respectively, as
background-like, and 16.5\% of events lacked a unanimous label. Differences
were especially clear during the intense 2019 sequence. On a common daily
grid, pairwise mean absolute differences between the smooth starting functions
reached 0.203 events per day, with correlations as low as 0.302. Nevertheless,
the optimiser reported successful termination for all three fits, the fitted
log-likelihoods differed by only 0.121, and the final background-rate
correlations all exceeded 0.999991. Only 0.38\%--1.17\% of the initial pairwise
mean absolute differences remained after fitting. A limited specification
analysis found little change in the triggering-parameter estimates. The
estimated background functions also remained similar to the baseline function
($r\geq0.983$), although their shapes were more affected by spline flexibility
than by the triggering-history length.

For this catalogue and baseline model, the three declustering-derived starting
functions led to essentially the same fitted ETAS solution. This provides
practical flexibility when selecting among the three tested procedures to
initialise a non-stationary ETAS background that is subsequently re-estimated.
However, it does not establish that other declustering algorithms, model
specifications, or arbitrary starting values would be interchangeable, nor
that declustering is unimportant when the event classifications themselves are
of scientific interest.
\end{abstract}

\section{Introduction}
\label{sec:introduction}

Earthquake catalogues record the times, locations, depths, and magnitudes of
detected seismic events. They provide the observational basis for estimating
seismicity rates, studying earthquake sequences, and constructing statistical
forecasting and hazard models. Two features make such catalogues statistically
challenging. First, earthquakes are not independent: one event can be followed
by an elevated rate of later events. Second, the recorded rate can change for
reasons other than earthquake-to-earthquake triggering, including changes in
the underlying seismic process and in the ability of a monitoring network to
detect small events. A useful model must therefore distinguish temporal
clustering from slower variation in the rate that remains after clustering has
been accounted for.

Catalogue completeness is an important preliminary consideration. The
magnitude of completeness, $M_c$, is the threshold above which events are
treated as sufficiently reliably detected for a specified region and time
period. If smaller events are missed, particularly during a dense sequence,
the apparent event rate and the inferred links between earthquakes can both be
distorted. Methods based on the Gutenberg--Richter frequency--magnitude
relationship assess whether the observed magnitude distribution has reached
the approximately exponential form expected above $M_c$
\citep{gutenberg1944frequency,wiemer2000minimum,woessner2005assessing}.
Comparative work shows that the performance of catalogue-based estimators can
vary when completeness is heterogeneous, so a single catalogue-wide threshold
should be treated as an operational modelling choice rather than a guarantee
of uniform detection through time \citep{wang2025magnitude}. This issue is
particularly relevant here because the study period includes the 2019
Ridgecrest sequence, whose magnitude 7.1 mainshock was followed by unusually
intense recorded activity \citep{ross2019hierarchical}.

Earthquake declustering attempts to separate events that appear to form
dependent sequences from events treated as background or as the principal
event of a sequence. The resulting labels are algorithmic classifications,
not directly observed causal categories. Common approaches embody different
definitions of proximity. Gardner--Knopoff uses magnitude-dependent spatial
and temporal windows \citep{gardner1974}; Reasenberg's approach constructs and
updates groups of potentially interacting events \citep{reasenberg1985}; and
nearest-neighbour methods use a combined time--space--magnitude distance to
identify unusually close event pairs \citep{zaliapin2013,zaliapin2020declustering}.
These distinctions matter because the algorithms need not identify the same
events as background-like.

Comparative studies have found that both the chosen algorithm and its parameter
settings can change the number and characteristics of events retained after
declustering \citep{perry2024comparative}. Such choices can also alter the
estimated Gutenberg--Richter $b$-value, which describes the relative frequency
of smaller and larger earthquakes \citep{mizrahi2021effect}. Recent work
comparing deterministic cluster structures with probabilities from an ETAS
model further emphasises that the two frameworks answer related but different
questions: deterministic procedures assign sequence labels, whereas ETAS
represents each event through model-based background and triggering
contributions \citep{spassiani2025reconciling}. Agreement between methods
should therefore not be assumed, and a high overall proportion of matching
labels can be misleading when one class is much more common than the other.

The epidemic-type aftershock sequence (ETAS) model addresses temporal
dependence by treating earthquake occurrence as a self-exciting point process
\citep{ogata1988statistical}. Its conditional event rate is the sum of a
background component and contributions from earlier earthquakes. Larger
events can produce larger triggering contributions, and each contribution
decays with elapsed time. In this statistical decomposition, ``background''
means the part of the fitted rate not attributed to the modelled influence of
earlier catalogue events; it should not be interpreted as an observed physical
status for an individual earthquake. A stationary temporal ETAS model assumes
that this background rate is constant. A non-stationary model relaxes that
assumption and allows the background rate to vary with time, for example by
representing it with splines and controlling excessive variation through a
roughness penalty \citep{kumazawa2014nonstationary,kattamanchi2017nonstationary}.

Declustering and ETAS can be connected in more than one way. ETAS itself can be
used for probabilistic declustering, with background and parent probabilities
estimated from the fitted branching structure \citep{zhuang2002stochastic}.
Alternatively, a conventional declustering calculation can provide a
preliminary estimate of the background rate before an ETAS model is fitted to
the complete event catalogue \citep{ogata1998space,zhuang2002stochastic}. In
the latter use, declustering does not determine which events enter the final
likelihood. It supplies initial parameter values for a numerical optimiser,
which then re-estimates both the background and triggering components from all
included events.

This distinction creates an optimisation question. ETAS likelihoods can be
difficult to optimise because background and triggering parameters are
interdependent, and numerical performance can depend on parameterisation and
starting values \citep{veen2008estimation,lombardi2015simulated}. Existing
work therefore establishes two separate concerns: alternative declustering
methods can produce different catalogues or event labels, and ETAS estimation
can be sensitive to numerical fitting choices. What remains less clear is
whether scientifically plausible differences between declustering methods
persist when they enter only as alternative initial background functions for
the same penalised non-stationary ETAS likelihood. Most declustering
comparisons study the resulting catalogues or compare deterministic labels
with ETAS probabilities; they do not isolate initialisation while holding the
events, likelihood, model structure, and all other starting values fixed.

This study addresses that narrower gap using the Southern California
Earthquake Data Center catalogue from 20 May 2000 to 20 May 2026
\citep{scedc2013,scedc_catalog_search2026}. The research question is: \emph{how
sensitive are fitted non-stationary temporal ETAS estimates to background-rate
initialisations derived from different declustering methods when the input
events and all other modelling choices are held constant?} Three contrasting
declustering schemes---Gardner--Knopoff, a deterministic nearest-neighbour
construction, and a Reasenberg-style interaction method---are compared at the
event-classification level. Their background-like events are then converted
into monthly rates and smooth spline starting functions. Finally, the same
penalised temporal ETAS model is fitted from each start, allowing the analysis
to trace differences from event labels, through the actual optimiser inputs,
to the final fitted background and triggering parameters. A limited
specification analysis examines whether the main parameter estimates also
change under moderate variations in spline flexibility and triggering-history
length.

The contribution is therefore methodological rather than a claim that one
declustering algorithm is correct. The controlled design distinguishes
disagreement in the original classifications from sensitivity of the final
model fit, while class-specific agreement measures and like-for-like curve
comparisons quantify where discrepancies arise and how much remains after
optimisation. Section~\ref{sec:methods} describes the catalogue,
completeness assessment, declustering implementations, construction of the
initial background functions, and ETAS model. Section~\ref{sec:results}
reports the classification, initialisation, fitted-model, and robustness
results. Section~\ref{sec:discussion} interprets the findings in relation to
previous declustering and non-stationary ETAS work and discusses their scope
and limitations.
